\documentclass[12pt]{article}
\usepackage{amsmath,amssymb}

\begin{document}
\section*{{2025 AIME II Problems/Problem 6}}
Source: AoPS Wiki

\subsection*{Solution}
Let $GH=2x$ and $GF=2y$ . Notice that since $\overline{BC}$ is perpendicular to $\overline{GH}$ (can be proven using basic angle chasing) and $\overline{BC}$ is an extension of a diameter of $\omega_1$ , then $\overline{CB}$ is the perpendicular bisector of $\overline{GH}$ . Similarly, since $\overline{AD}$ is perpendicular to $\overline{GF}$ (also provable using basic angle chasing) and $\overline{AD}$ is part of a diameter of $\omega_1$ , then $\overline{AD}$ is the perpendicular bisector of $\overline{GF}$ . From the Pythagorean Theorem on $\triangle GFH$ , we have $(2x)^2+(2y)^2=12^2$ , so $x^2+y^2=36$ . To find our second equation for our system, we utilize the triangles given. Let $I=\overline{GH}\cap\overline{CB}$ . Then we know that $GFBI$ is also a rectangle since all of its angles can be shown to be right using basic angle chasing, so $FG=IB$ . We also know that $CI+IB=2\cdot 15=30$ . $IA=y$ and $AB=6$ , so $CI=30-y-6=24-y$ . Notice that $CI$ is a height of $\triangle CHG$ , so its area is $\frac{1}{2}(2x)(24-y)=x(24-y)$ . Next, extend $\overline{DA}$ past $A$ to intersect $\omega_2$ again at $D'$ . Since $\overline{BC}$ is given to be a diameter of $\omega_2$ and $\overline{BC}\perp\overline{AD}$ , then $\overline{BC}$ is the perpendicular bisector of $\overline{DD'}$ ; thus $DA=D'A$ . By Power of a Point, we know that $CA\cdot AB=DA\cdot AD'$ . $CA=30-6=24$ and $AB=6$ , so $DA\cdot AD'=(DA)^2=24\cdot6=144$ and $DA=D'A=12$ . Denote $J=\overline{DA}\cap\overline{GF}$ . We know that $DJ=DA-AJ=12-x$ (recall that $GI=IH=x$ , and it can be shown that $GIAJ$ is a rectangle). $\overline{DJ}$ is the height of $\triangle DGF$ , so its area is $\frac{1}{2}(2y)(12-x)=y(12-x)$ . We are given that $[DGF]=[CHG]$ ( $[ABC]$ denotes the area of figure $ABC$ ). As a result, $x(24-y)=y(12-x)$ . This can be simplified to $y=2x$ . Substituting this into the Pythagorean equation yields $5x^2=36$ and $x=\frac{6}{\sqrt{5}}$ . Then $y=\frac{12}{\sqrt{5}}$ . $[EFGH]=2x\cdot2y=2\cdot\frac{6}{\sqrt{5}}\cdot2\cdot\frac{12}{\sqrt{5}}=\frac{288}{5}$ , so the answer is $288+5=\boxed{293}$ . ~ eevee9406 ~ Edited by aoum

\end{document}
